\documentclass[12pt]{article}

%%\usepackage[english]{babel}
\usepackage{amsmath}
\usepackage{amssymb}
\usepackage{url}


\begin{document}

\title{Econometrics Notes VII: Elementary Statistics - Part 4: }
\author{Artur Wegrzyn}
\maketitle

\abstract{This note presents basic results concerning distributions that will show up frequently in later considerations: exponential, gamma and $\chi^2$. }

\tableofcontents


\subsection{Fisher's Theorem}
In this and subsequent section use is made of the Fisher's theorem that states 
that given the sample discussed, the sample mean and sample variance are independent 
random variables. I.e. given:
\begin{equation*}
\overline{X} = \frac{1}{n} \sum_{i=1}^{n} X_i
\end{equation*}

and:
\begin{equation*}
\tilde{S}^2 = \frac{1}{n-1} \sum_{i=1}^{n} (X_i - \overline{X})^2
\end{equation*}

we have:
\begin{equation*}
\overline{X} \perp\!\!\!\perp \tilde{S}^2
\end{equation*}

Proof of this theorem will be given in a separate note. For now, note that the symbol 
$\perp\!\!\!\perp$ is commonly used to denote independence of random variables. Also, 
recall that in case of normal distriubtion independence and orthogonality 
are equivalent - this will be a topic of yet separate note.


\section{Gamma Distribution}
Let $Y_1, Y_2, \dots, Y_n$ be a sample drawn indepentently from an exponential 
distribution with expected value $1 / \lambda$:

\section{Square of Standard Normal Distribution}
Suppose that $U$ follows standard normal distribution:
\begin{equation*}
U \sim N(0, 1)
\end{equation*}

Consider $Z = U^2$. Then:
\begin{equation*}
P(Z \leq z) = P(U^2 \leq z) = P(-\sqrt{z} < U \leq \sqrt{z})
\end{equation*}

Hence:
\begin{equation*}
f_{Z}(z) = \frac{dF_{Z}(z)}{dz} = \frac{d}{dz} P(-\sqrt{z} < U \leq \sqrt{z})
\end{equation*}

A little bit of algebra yields:
\begin{equation*}
f_{Z}(z) = 
\end{equation*}



\subsection{Statistical Motivation}
Given the sample, consider the biased variance estimator:
\begin{equation*}
S_{XY} = \frac{1}{n} \sum_{i=1}^{n} (X_i - \overline{X})^2 
\end{equation*}

The aim is to find its distribution. As previously proved, this estimator can be rewritten as:
\begin{equation*}
S_{XY} = \frac{1}{n} \sum_{i=1}^{n} (X_i - \mu)^2 - (\overline{X} - \mu)^2
\end{equation*}


\subsection{Derivation of PDF}

\section{Derivation of t-Student Distribution PDF}
\subsection{t-Student Distribution - Statistical Motivation}

\subsection{Derivation of the PDF}

\section{Derivation of F-Snedecor Distribution PDF}
\subsection{F-Snedecor Distribution - Statistical Motivation}
\subsection{Derivation of the PDF}

\bibliography{tsa_bibliography}
\bibliographystyle{ieeetr}

\end{document}
